\chapter{結論}

\section{本研究の貢献}

本研究の主要な貢献は以下の3点である。

\textbf{（1）地形考慮型経路計画手法の提案：}
RTAB-Mapから得られる点群データを用い、傾斜角、標高差、表面粗さの3要素を統合した地形複雑度指標を構築した。これをA*の評価関数に組み込むことで、地形を考慮した経路計画を実現した。従来手法が経路長のみを最適化していたのに対し、本手法は地形の走行適性も同時に評価することで、より実用的な経路生成を可能にした。

\textbf{（2）地形考慮型経路計画の2つの実装形態の提案：}
グリッドベースの探索に地形コストを統合したTerrain-Aware A*と、補間技術により経路平滑化と地形考慮を両立するField D* Hybridを開発し、96シナリオでの包括的な比較実験により性能を定量的に評価した。実験結果から、TA*は最短経路（21.27m）と96.88\%の成功率を達成し、効果量Cohen's $d=0.91$で統計的に有意な性能改善を示した。Field D* Hybridは計算時間0.175秒で100\%の成功率を達成し、実時間性と信頼性を両立した。これにより、ロボットの用途や環境特性に応じた手法選択の指針を提供した。

\textbf{（3）複数評価指標による多角的評価の実施：}
地形コストの評価に点平均法とセグメント加重法の2つの方法を導入し、アルゴリズムの地形回避能力と実用的な移動コストを多角的に評価した。さらに、箱ひげ図、地形複雑度別分析、失敗シナリオ分布など詳細な分析により、各手法の性能特性を明確化した。点平均法はアルゴリズムの設計意図通りの動作を検証する指標として、セグメント加重法は実用的な移動コストを評価する指標として、それぞれ異なる役割を果たすことを明確にした。

これらの成果により、不整地環境における自律移動ロボットの安全かつ効率的な経路計画の実現に貢献した。

\section{結言}

本研究では、不整地環境における自律移動ロボットの地形考慮型経路計画フレームワークを構築し、大規模ベンチマークによる実験と統計的検証により、その有効性を実証した。

\textbf{主要な成果：}
第一に、標高差、傾斜角、表面粗さを統合した地形複雑度指標を定式化し、予備実験により最適パラメータ（$\alpha=\beta=0.4$, $\gamma=0.2$, $w_t=1.0$）を決定した。RTAB-Mapとの統合により実装可能性を示し、0.2m解像度での実時間地形評価を実現した。

第二に、96シナリオ（3種のマップサイズ$\times$3種の地形複雑度）の大規模ベンチマークにより、4手法の性能を統計的に検証した。TA*は最短経路21.27m、成功率96.88\%（93/96）、地形コスト削減40.3\%を達成し、AHA*との比較でWelch t-test $p<0.001$、Cohen's $d=0.91$の大きな効果量により統計的に有意な改善を実証した。Field D* Hybridは計算時間0.175秒で100\%の成功率を達成し、実時間性（0.175秒）と高信頼性（100\%）を両立した最もバランスの取れた手法であることが確認された。これは、実環境への即座の実装を可能にする実用レベルの性能である。

第三に、箱ひげ図、地形複雑度別分析、失敗シナリオ分布など詳細な分析により、各手法の性能特性を定量的に解明した。TA*は低・中複雑度環境（<0.55）で100\%の成功率を達成し、高複雑度環境（$\geq$0.55）で90.6\%を示した。失敗要因分析により、全ての失敗が地形複雑度0.65以上の超高複雑度環境でタイムアウト（300秒）に起因することを特定し、実装上の注意点を明確化した。

\textbf{実用化への貢献：}
本フレームワークは、農業ロボット（畝回避、経路効率化）、災害対応ロボット（瓦礫回避、迅速な救助活動）、惑星探査ローバー（未知地形での自律探査）など、多様な分野への応用が可能である。特に、Field D* Hybridの実時間性（0.175秒）は、実環境でのリアルタイム経路計画を実現する。

96シナリオでの包括的評価により、地形考慮型経路計画が実用レベルに到達したことを実証した。今後の課題として、(1) 実機ロボットによるフィールド実証実験、(2) 動的環境への対応、(3) 機械学習による地形予測の統合が挙げられる。本研究が、不整地環境における自律移動ロボットの安全性と効率性の大幅な向上に寄与し、フィールドロボティクスの更なる発展に貢献することを期待する。
